\section{Finite estimates and the irreducible law}\label{c:inputs}
We first identify the inputs that do not depend on prescribed endpoints.
A strip of integral height $h$ has physical thickness $dh$, where
$d=\sqrt3/2$. Write $B_h$ for its critical bridge mass from one fixed
bottom port, summed over all compatible top ports, and put $B_0=1$.
A strict bridge stays between its boundary cuts. Its terminal offset lies
on the row-staggered lattice $\mathbb Z+h/2$ in horizontal side units.
There is no translation factor in $B_h$.

\begin{proposition}[Finite boundary and length inputs]\label{c:finite}
The following estimates hold in the port convention above.
\begin{enumerate}
\item For every integer $h\ge0$, $B_h\asymp(1+h)^{-1/4}$.
For all sufficiently large integers $h$, the length-weighted
bridge mass is at most $Ch^{13/12}$, and
$B_h\{L\ge ch^{4/3}\}\ge ch^{-1/4}$.
For every integer $h\ge1$ and every $r\ge0$, the mass of
bridges of diameter exceeding $r$ is at most $C(1+r)^{-1/4}$.
\item In a convex pure domain, the mass from one boundary source of
completed paths of diameter at least $r$ is at most $C(1+r)^{-1/4}$
for $r\ge0$.
For the half-plane arch kernel and integer $l\ge1$,
$K_l\asymp l^{-5/4}$ and $K_{l+1}\le K_l$.
The same monotonicity comparison is valid along one straight side
of a convex pure truncation whenever both compared adjacent ports
remain on that side. Arches of diameter at most
$Al$ retain mass $cl^{-5/4}$ for a fixed $A$.
\item From a prescribed port $a$ with inward pure normal $n$, consider
the terminal line $m\cdot(z-a)=s$, where $m$ is its outward pure unit
normal and the angle between $n$ and $m$ is at most $\pi/3$.
Here $s>0$ is the Euclidean distance from $a$ to the terminal line.
Crossings that stay in the initial half-plane and before the terminal
line until their final exit have mass comparable to $s^{-1/4}$.
For any fixed positive relative tolerance, and all sufficiently large
admissible $s$ with a threshold depending only on that tolerance, the lower
bound can be confined within distance equal to that tolerance times $s$
of the segment $a+[0,s]m$. The corresponding two-cut upper
bound holds without confinement.
\item In a convex pure domain of diameter $O(R)$, the total
length-weighted mass of chords, summed over ordered boundary endpoint
pairs, is at most $R^{25/12+o(1)}$.
\end{enumerate}
\end{proposition}
These are the finite results of \cite[Theorem 1.1, Propositions 7.1, 7.4, 10.1 and Corollary 8.1]{UniformMarked}. The statement records the input
before use; the full analytic and finite geometric proofs are in that
companion. In particular its free-boundary mean is a finite input, not a
consequence of the endpoint theorem here. The disk-transfer construction
\cite[Theorem 1.1]{DiskTransfer} provides an independent fixed-corridor first moment.
We use the specified finite method at each application below; in particular
the logarithmic confinement in the polynomial-vacuum route retains its
own precision.

We use also the local contour identity with phase
$e^{i\sigma W}$, $\sigma=3/8$, where $W$ is total turning.
This is the phase $e^{-5iW/8}$ of the usual observable multiplied by
the terminal tangent relative to the initial direction;
see \cite[Section 2, Lemma 1]{DuminilCopinSmirnov2012}.
At an unvisited center the two outgoing weights sum to the incoming
weight, since $2\kappa\cos(\pi/8)=1$. Prefixes returning to a visited
center cancel in pairs under reversal of their unfinished simple loop:
their turn increments differ by $8\pi/3$, and hence their phases have
opposite signs. Summing over a finite domain gives signed exit mass one.
For a convex domain the real exit coefficients lie in
$[\cos(3\pi/8),1]$. Subtracting the identity in a convex subdomain
bounds the mass of paths lost on restriction by a fixed multiple of the
mass to the newly exposed sides. Exhaustion is justified by the diameter
bound in Proposition~\ref{c:finite}. This is the local winding mechanism
of \cite{DuminilCopinSmirnov2012}; below we explain the additional arc
and spectator constructions when they are used.

An irreducible strict bridge has positive integral height and no
intermediate singly crossed level. Every bridge has a unique ordered
factorization into such pieces. The honeycomb mid-edge version of the
irreducible construction is given in \cite[Appendix]{BeatonEtAl2014}.
Let $p(\gamma)=\kappa^{L(\gamma)}$
on irreducibles from a fixed port. Let $H,Z,L,D$ denote respectively the
height, transverse displacement, length and diameter of one piece.

\begin{proposition}[Renewal inputs]\label{c:renewal}
The mass $p$ is a probability. Successive complete pieces have independent
law $p$, and, for $s>0$,
\[
 \sum_{h\ge0}B_he^{-sh}=
 (1-\mathbb E_p e^{-sH})^{-1}\asymp s^{-3/4}\quad(s\downarrow0).
\]
Moreover,
\[
 p(H>x)+p(D>x)\le Cx^{-3/4},\quad
 \mathbb E_p[L;H\le x]\le x^{7/12+o(1)},
\]
\[
 p(L>n)\le n^{-9/16+o(1)},\qquad
 1-\mathbb E_p e^{-uL}=u^{9/16+o(1)}.
\]
Almost surely, for independent pieces, their first $k$ heights and lengths
sum respectively to $k^{4/3+o(1)}$ and $k^{16/9+o(1)}$, and their
diameters sum to at most $k^{4/3+o(1)}$.
\end{proposition}
The proofs and the full joint-law conventions are
\cite[Propositions 2.2, 2.4, Remark 2.5 and Lemma 2.6]{Renewal}. Only the independent \emph{pieces}
are asserted here; their coordinates are not independent. $B_h$ is a
renewal occurrence mass, not a probability conditioned on a future height.
No estimate in this proposition conditions a piece on its length.

Throughout, $f(R)=R^{a+o(1)}$ means
$R^{a-\varepsilon}\le f(R)\le R^{a+\varepsilon}$ for every fixed
$\varepsilon>0$ at all sufficiently large scales, with fixed geometric
parameters chosen first. A one-sided use retains just the corresponding
inequality. All suprema of atoms are on their actual support lattice;
off-lattice probabilities are zero.
